\section{Discussion}
\label{sec:discussion}

The central observation of this paper is that stereographic
projection, composed with standard QSP, converts the bounded
polynomial output of the circuit into an unbounded rational
function---without modifying the quantum algorithm itself.  The
basis change identity
$S_z = V^{-1}W(\rtilde)V$ with $V = \mathrm{diag}(1,i)$
(Proposition~\ref{prop:basis-change}) makes this reduction
precise: the full machinery of standard QSP---phase-finding
algorithms~\cite{DongMengWhaleyEtAl2021,Haah2019}, the
function class characterization, and convergence theory---applies
immediately, and the ``unboundedness'' is a free consequence of
the encoding/decoding interface.  The stereographic viewpoint
also illuminates standard QSP itself: even the standard
protocol, when decoded via stereographic projection, produces
unbounded rational functions of the real-line parameter
(Section~\ref{subsec:further-consequences}).

As discussed in Section~\ref{sec:applications}, the stereographic
variable substitution also provides a complementary approach to
QSVT-based quantum state
preparation~\cite{McArdleGilyenBerta2025}, replacing the arcsine
encoding with the stereographic map and connecting the
approximation theory to Boyd's rational Chebyshev
basis~\cite{Boyd1990,Boyd1987}.  While state preparation
amplitudes remain bounded---one cannot encode an unbounded function
directly into a quantum state---the stereographic basis is
particularly well suited to functions defined on unbounded domains
or with naturally rational structure.

At the same time, the construction has clear limitations.
The measurement overhead grows as
$\mathcal{O}(r^6/\epsilon^2)$ for large eigenvalues
(Eq.~\eqref{eq:shot-count}), making the approach impractical
when $r \gg 1$.  Lifting the single-qubit theory to a
multi-qubit eigenvalue transformation is straightforward when
the eigenbasis of the Hamiltonian is known, but this is a
strong assumption that restricts applicability to structured
operators.  The practical advantage window is depth-limited
devices processing operators with moderate eigenvalues and
known eigenbasis structure; a full development of the
multi-qubit framework is presented in a companion
paper~\cite{KharaziLeytonOrtega2026b}.

Our construction rests on the QSP/QSVT framework established
by Low and Chuang~\cite{LowChuang2017,LowChuang2019} and
Gily\'en et al.~\cite{GSLW2019}, with comprehensive reviews by
Martyn et al.~\cite{MartynRossiTanChuang2021} and
Lin~\cite{Lin2022}.  The formal separation between bounded and
unbounded polynomial approximations established by Tang and
Tian~\cite{TangTian2024} provides the theoretical motivation:
it shows that the boundedness constraint is a genuine barrier,
not merely a technical inconvenience.  Among extensions,
GQSP~\cite{MotlaghWiebe2023} is the closest in spirit: it
lifts the parity and realness constraints on~$P$ and provides
dramatically faster angle-finding, but retains the boundedness
constraint $|P| \leq 1$.  Stereographic QSP is
complementary---it lifts boundedness but not parity, and
produces rational rather than polynomial output.  Quantum Phase
Processing~\cite{WangZhangYuWang2023}, modular
QSP~\cite{RossiCeroniChuang2025}, Randomized
QSVT~\cite{WangZhangHazraEtAl2025}, Hamiltonian block
encoding~\cite{KangSu2025}, and GQSP-based matrix function
synthesis~\cite{MahasingheEtAl2025} each address different
costs of the standard framework while retaining bounded output;
our work changes the classical encoding/decoding interface
rather than the quantum algorithm.  The connection to Boyd's
rational Chebyshev
functions~\cite{Boyd1990,Boyd1987} links stereographic QSP to
classical approximation theory on unbounded domains, suggesting
that convergence rates for rational Chebyshev expansions
translate directly to QSP performance guarantees.

The framework opens several directions for further work.  On
the algorithmic side, finding QSP phases for a given target
rational function $f(r)$ requires decomposing $f$ into a
polynomial pair $(P, Q)$ satisfying the normalization
constraint before invoking standard phase-finding; GQSP's
efficient angle-finding~\cite{MotlaghWiebe2023} may simplify
this step.  On the theoretical side, the numerical results in
Figure~\ref{fig:target-gallery} demonstrate convergence for
several targets, with the inset in panel~(e) showing that
eigenvalue inversion converges fastest while bounded targets
converge more slowly---understanding the precise relationship
between rational Chebyshev convergence rates and the QSP
normalization constraint is an open question.  The natural
error metric on unbounded domains is the weighted $L^2$ norm
with Cauchy--Lorentz weight $1/(1+r^2)$, which downweights
large-$r$ errors where stereographic compression is most
severe.

On the multi-qubit side, the key open question is whether
the eigenbasis requirement of stereographic block
encoding~\cite{KharaziLeytonOrtega2026b} can be
relaxed---perhaps by combining stereographic encoding with
quantum phase estimation, or by constructing a stereographic
analogue of the QSVT product for block-encoded matrices.
Extending the class of tractable Hamiltonians to sparse models
with known eigenbasis structure, integrable models, or
tensor-network states would establish the practical scope.  At a more fundamental level,
replacing the phase rotations $\exp[i\phi\PauliZ]$ with general
signal-type operators $O_\alpha$ parameterized by
$\alpha \in \C\mathbb{P}^1$ yields a generalized normalization
condition
\begin{equation}
	|P|^2 + |Q|^2 = (1{+}r^2)^d \textstyle\prod_j(1{+}|\alpha_j|^2),
	\label{eq:generalized-norm}
\end{equation}
whose forward direction follows from the same inductive
argument as Theorem~\ref{thm:unnormalized}, but whose
converse---decomposing arbitrary polynomial pairs into such a
product---remains open.

The stereographic projection is one of the oldest maps in
mathematics, but its application to quantum signal processing
appears to be new.  The fact that a simple change of
coordinates can bypass a fundamental constraint of QSP suggests
that the interface between encoding and processing---the way
classical data enters and exits the quantum circuit---deserves
systematic study.  Other geometric constructions (hyperbolic
encodings, projective coordinates) may yield analogous
extensions, and the general question of which function classes
are realizable under different encoding/decoding strategies
remains open.
